\begingroup
\color{borrador2}
\textcolor{borrador2}{La siguiente figura presenta una alternativa del recorrido completo que incorpora la decisión de ámbito y los eventos transmitidos por la aplicación web.}

\begin{figure}[H]
  \centering
  \resizebox{0.97\textwidth}{!}{%
    \begin{tikzpicture}[
      proceso/.style={draw=borrador2, thick, rounded corners=3pt, fill=white,
                      text=borrador2, align=center, font=\small,
                      minimum width=3.0cm, minimum height=1.0cm},
      decision/.style={diamond, draw=borrador2, thick, fill=borrador2!7,
                       text=borrador2, align=center, font=\footnotesize,
                       aspect=2.1, inner sep=2pt},
      salida/.style={proceso, fill=borrador2!9},
      flecha/.style={-{Latex[length=2.2mm]}, thick, draw=borrador2},
      etiqueta/.style={font=\scriptsize, text=borrador2, fill=white, inner sep=1pt},
    ]
      \node[proceso] (pregunta) at (0,0) {Recibir pregunta};
      \node[decision] (fija) at (3.7,0) {¿Tiene\\respuesta fija?};
      \node[proceso] (ambito) at (7.4,0) {Decidir y validar\\el ámbito};
      \node[proceso] (recuperar) at (11.1,0) {Buscar y aplicar\\los tres cortes};
      \node[decision] (fragmentos) at (14.8,0) {¿Quedan\\fragmentos?};

      \node[salida] (cerrar) at (0,-3.1) {Anotar, registrar\\y sugerir};
      \node[salida] (final) at (3.7,-3.1) {Respuesta final};
      \node[decision] (inventada) at (7.4,-3.1) {¿Nombra una\\titulación ajena?};
      \node[proceso] (generar) at (11.1,-3.1) {Emitir fuentes y\\transmitir generación};
      \node[salida] (ausencia) at (14.8,-3.1) {Ausencia de\\información};

      \node[salida] (inmediata) at (3.7,-5.5) {Respuesta inmediata};
      \node[salida] (retirar) at (7.4,-5.5) {Emitir \texttt{borrar} y\\sustituir el texto};

      \begin{scope}[on background layer]
      \draw[flecha] (pregunta) -- (fija);
      \draw[flecha] (fija) -- node[etiqueta]{no} (ambito);
      \draw[flecha] (ambito) -- (recuperar);
      \draw[flecha] (recuperar) -- (fragmentos);
      \draw[flecha] (fragmentos) -- node[etiqueta, right]{no} (ausencia);
      \draw[flecha] (fragmentos) -- node[etiqueta, sloped]{sí} (generar);
      \draw[flecha] (generar) -- (inventada);
      \draw[flecha] (inventada) -- node[etiqueta]{no} (final);
      \draw[flecha] (final) -- (cerrar);
      \draw[flecha] (fija) -- node[etiqueta, right]{sí} (inmediata);
      \draw[flecha] (inventada) -- node[etiqueta, right]{sí} (retirar);
      \draw[flecha] (inmediata.west) -| (cerrar.south);
      \draw[flecha] (retirar.north) -- ++(0,0.7) -| (final.south);
      \draw[flecha] (ausencia.south) -- ++(0,-3.8) -| (cerrar.south);
      \end{scope}
    \end{tikzpicture}%
  }
  \caption[Alternativa del recorrido de una consulta]{\textcolor{borrador2}{Versión alternativa del recorrido de una consulta. Distingue la decisión de ámbito de la generación, mantiene las salidas sin generación cuando corresponden y añade las fuentes, la retirada de texto, el registro y las sugerencias del flujo web. Fuente: Elaboración propia.}}
  \label{fig:secuencia_rag_alt}
\end{figure}
\endgroup
